
%\usepackage[extreme]{savetrees}

\usepackage[a4paper]{geometry}
\usepackage{pgfplots}
\pgfplotsset{compat=1.9}

\usepgfplotslibrary{colorbrewer}
\pgfplotsset{cycle list/Blues-9}

\usepackage{parskip} %swap indentation to spacing


\usepackage[justification=centering]{caption}
\usepackage{subcaption}




\usepackage{graphicx}
\usepackage{float}
\usepackage{amsmath}
\usepackage{booktabs}
\usepackage{tabularx}
\usepackage{perpage}
\usepackage{smartdiagram}
\usepackage{cancel}
\usepackage{amssymb}
\usepackage{bm}
\usepackage{longtable}
%\usepackage{minted}
\usepackage{listings}
\usepackage{courier}
\usepackage{fancybox}

\newcolumntype{M}[1]{>{\centering\arraybackslash}m{#1}}
\newcolumntype{L}[1]{>{\raggedright\arraybackslash}p{#1}}
\newcolumntype{R}[1]{>{\raggedleft\arraybackslash}p{#1}}

\usepackage{hhline}

\usepackage{bm}
\usepackage{siunitx}
\usepackage[utf8]{inputenc}
\usepackage[english]{babel}
\usepackage{graphicx}
\usepackage{amsthm}

\usepackage{polynom}

\usepackage{tabularray}

\usepackage{verbatim}

\usepackage{array}
\usepackage{ulem}
\usepackage{mathtools}
\usepackage{amsfonts}

\usepackage{algorithm}
\usepackage{algpseudocode}


\usepackage{color}   %May be necessary if you want to color links
\usepackage[pagebackref=true]{hyperref}
\hypersetup{
    colorlinks=true, %set true if you want colored links
    linktoc=all,     %set to all if you want both sections and subsections linked
    linkcolor=black, %choose some color if you want links to stand out
}

\addto\extrasenglish{% comment out if you want autorefs to say "subsection" instead of using the symbol
    \def\chapterautorefname{\S\!}
    \def\sectionautorefname{\S\!}
    \def\subsectionautorefname{\S\!}
    \def\subsubsectionautorefname{\S\!}
}

\usepackage{cleveref}


\crefformat{chapter}{\S#2#1#3}
\crefmultiformat{chapter}{\S\S#2#1#3}{ and~#2#1#3}{, #2#1#3}{, and~#2#1#3}

\crefformat{section}{\S#2#1#3}
\crefmultiformat{section}{\S\S#2#1#3}{ and~#2#1#3}{, #2#1#3}{, and~#2#1#3}


\usepackage[explicit]{titlesec}

\setcounter{tocdepth}{3}
\setcounter{secnumdepth}{3}
\titleformat{\chapter}[display]
    {\normalfont\huge\bfseries}{\chaptertitlename\ {\fontfamily{cmr}\selectfont\thechapter}}{20pt}{\hyperlink{chap-\thechapter}{\Huge#1}
\addtocontents{toc}{\protect\hypertarget{chap-\thechapter}{}}}
\titleformat{name=\chapter,numberless}
    {\normalfont\huge\bfseries}{}{-20pt}{\Huge#1}
\titleformat{\section}
    {\normalfont\Large\bfseries}{\thesection}{1em}{\hyperlink{sec-\thesection}{#1}
\addtocontents{toc}{\protect\hypertarget{sec-\thesection}{}}}
\titleformat{name=\section,numberless}
    {\normalfont\Large\bfseries}{}{0pt}{#1}
\titleformat{\subsection}
    {\normalfont\large\bfseries}{\thesubsection}{1em}{\hyperlink{subsec-\thesubsection}{#1}
\addtocontents{toc}{\protect\hypertarget{subsec-\thesubsection}{}}}
\titleformat{name=\subsection,numberless}
    {\normalfont\large\bfseries}{\thesubsection}{0pt}{#1}

\titleformat{\subsubsection}
    {\normalfont\bfseries}{\thesubsubsection}{1em}{\hyperlink{subsubsec-\thesubsubsection}{#1}
\addtocontents{toc}{\protect\hypertarget{subsubsec-\thesubsubsection}{}}}
\titleformat{name=\subsubsection,numberless}
    {\normalfont\bfseries}{\thesubsubsection}{0pt}{#1}


\usepackage{xparse}
\usepackage{tikz}
\usetikzlibrary{calc}
\usetikzlibrary{arrows}
\usetikzlibrary{arrows.meta}
\usetikzlibrary{decorations.pathreplacing}
\usetikzlibrary{decorations.markings}
\usetikzlibrary{backgrounds}
\usetikzlibrary{positioning}
\usetikzlibrary{intersections}
\usetikzlibrary{calligraphy}
\usetikzlibrary{patterns}
\usepgflibrary{shadings}


\tikzset{
  show curve controls/.style={
    decoration={
      show path construction,
      curveto code={
        \draw[#1!50]
        (\tikzinputsegmentfirst)
        -- (\tikzinputsegmentsupporta)
        -- (\tikzinputsegmentsupportb)
        -- (\tikzinputsegmentlast)
        ;
        \fill[#1!50] (\tikzinputsegmentsupporta) circle(1pt);
        \fill[#1!50] (\tikzinputsegmentsupportb) circle(1pt);
        \draw[#1,line width=1pt]
        (\tikzinputsegmentfirst)
        .. controls (\tikzinputsegmentsupporta)
                and (\tikzinputsegmentsupportb) ..
        (\tikzinputsegmentlast);
      }
    },decorate
  }
}

\usepackage{tikz-cd}
\usepackage{quiver}

\tikzset{Arrow Style/.style={text=black, font=\boldmath}}

\newcommand{\tikzmark}[1]{%
    \tikz[overlay, remember picture, baseline] \node (#1) {};%
}

\NewDocumentCommand{\DrawArrow}{s O{} m m m}{%
    \begin{tikzpicture}[overlay,remember picture]
        \draw[->, thick, Arrow Style, #2] 
                ($(#3.west)+(\XShift,\YShift)$) -- 
                ($(#4.east)+(-\XShift,\YShift)$)
        node [midway,above] {#5};
    \end{tikzpicture}%
}



\newcommand\eqset{\mathrel{\overset{\makebox[0pt]{\mbox{\normalfont\tiny\sffamily set}}}{=}}}
\newcommand\eqdef{\mathrel{\overset{\makebox[0pt]{\mbox{\normalfont\tiny\sffamily def}}}{=}}}
\newcommand\eqquest{\mathrel{\overset{\makebox[0pt]{\mbox{\normalfont\tiny\sffamily ?}}}{=}}}

\newcommand{\triangleleftneq}{\mkern-1mu\mathrel{\ooalign{\adjustbox{scale={1}{0.575}, raise={0.475ex}}{\hspace*{0.55em}$\mid$}\cr$\lneq$\cr}\mkern-1mu}}
\newcommand{\trianglerightneq}{\mkern-1mu\mathrel{\ooalign{\adjustbox{scale={1}{0.575}, raise={0.475ex}}{\hspace*{-0.05em}$\mid$}\cr$\gneq$\cr}\mkern-1mu}}


\newcommand{\ihat}{\boldsymbol{\hat{\textbf{\i}}}}
\newcommand{\jhat}{\boldsymbol{\hat{\textbf{\j}}}}


\newcommand{\image}{\operatorname{im}}
\newcommand{\nullity}{\operatorname{null}}
\newcommand{\rank}{\operatorname{rank}}

\newcommand{\proj}{\operatorname{proj}}
\newcommand{\coker}{\operatorname{coker}}

\newcommand{\type}{\operatorname{type}}

\newcommand{\Sym}{\operatorname{Sym}}
\newcommand{\Alt}{\operatorname{Alt}}
\newcommand{\Aut}{\operatorname{Aut}}
\newcommand{\End}{\operatorname{End}}
\newcommand{\Orb}{\operatorname{Orb}}
\newcommand{\Stab}{\operatorname{Stab}}
\newcommand{\Conj}{\operatorname{Cl}}
\newcommand{\fix}{\operatorname{fix}}
\newcommand{\Syl}{\operatorname{Syl}}
\newcommand{\lcm}{\operatorname{lcm}}



\makeatletter
\newcommand{\pgfplotsdrawaxis}{\pgfplots@draw@axis}
\makeatother
\pgfplotsset{only axis on top/.style={axis on top=false, after end axis/.code={
             \pgfplotsset{axis line style=opaque, ticklabel style=opaque, tick style=opaque,
                          grid=none}\pgfplotsdrawaxis}}}

                          
\newcommand{\drawle}{-- (rel axis cs:1,0) -- (rel axis cs:1,1) -- (rel axis cs:0,1) \closedcycle}
\newcommand{\drawge}{-- (rel axis cs:1,1) -- (rel axis cs:1,0) -- (rel axis cs:0,0) \closedcycle}



\makeatletter
\renewcommand*\env@matrix[1][*\c@MaxMatrixCols c]{%
  \hskip -\arraycolsep
  \let\@ifnextchar\new@ifnextchar
  \array{#1}}
\makeatother

\newcommand{\Mod}[1]{\ (\mathrm{mod}\ #1)}

\newcommand*{\XShift}{0.5em}
\newcommand*{\YShift}{0.5ex}
\definecolor{GraphRed}{HTML}{ff0000}
\definecolor{GraphBlue}{HTML}{00b0f0}
\definecolor{GraphPink}{HTML}{fa23ff}
\definecolor{GraphGreen}{HTML}{00b050}
\definecolor{GraphOrange}{HTML}{ffc000}
\definecolor{light-gray}{gray}{0.95}


\allowdisplaybreaks[4]

\makeatletter
\gdef\@ctrerr{%
  \@latex@warning{Counter too large}}
\makeatother
% changes counter too large errors into warnings.
% useful for documents with a large number of footnotes








\renewcommand{\thefootnote}{\fnsymbol{footnote}}
\MakePerPage{footnote}

\newtheorem{theorem}{Theorem}[section]
\newtheorem*{theorem*}{Theorem}
\newtheorem{corollary}{Corollary}[theorem]
\newtheorem*{corollary*}{Corollary}
\newtheorem{lemma}[theorem]{Lemma}
\newtheorem*{lemma*}{Lemma}
\newtheorem{proposition}[theorem]{Proposition}
\newtheorem*{proposition*}{Proposition}

\newtheorem*{claim}{Claim}

\newtheorem{conjecture}[theorem]{Conjecture}
\newtheorem*{conjecture*}{Conjecture}

\theoremstyle{definition}
\newtheorem{definition}{Definition}[section]

\theoremstyle{remark}
\newtheorem*{remark}{Remark}
\newtheorem*{exercise}{Exercise}
\newtheorem*{exercises}{Exercises}

\newtheorem*{example}{Example}
\AtBeginEnvironment{example}{%
  \pushQED{\qed}\renewcommand{\qedsymbol}{$\triangle$}%
}
\AtEndEnvironment{example}{\popQED\endexample}

\renewcommand\qedsymbol{$\blacksquare$}







\DeclareMathOperator{\sech}{sech}
\DeclareMathOperator{\csch}{csch}
\DeclareMathOperator{\arcsec}{arcsec}
\DeclareMathOperator{\arccot}{arccot}
\DeclareMathOperator{\arccsc}{arccsc}
\DeclareMathOperator{\arcosh}{arcosh}
\DeclareMathOperator{\arsinh}{arsinh}
\DeclareMathOperator{\artanh}{artanh}
\DeclareMathOperator{\arsech}{arsech}
\DeclareMathOperator{\arcsch}{arcsch}
\DeclareMathOperator{\arcoth}{arcoth} 

\usepackage[most]{tcolorbox}

\newtcolorbox{tcbsiderules}[1][]{blanker, breakable, 
     left=3mm, right=3mm, top=1mm, bottom=1mm,
     borderline west={1pt}{-1.5pt}{black},
     borderline west={0.5pt}{0pt}{black},
     borderline east={0.5pt}{0pt}{black},
     borderline east={1pt}{-1.5pt}{black},
     before upper=\indent, parbox=false, #1}

\newtcolorbox{axiombox}[1][]{
    left=5pt,right=5pt,
    colback=cyan!10!white,
    frame empty,
    width=0.98\textwidth,
    center,
    #1
}


\usepackage[myheadings]{fullpage}
\usepackage{fancyhdr}
\usepackage{lastpage}
\usepackage{float}
\usepackage{url, lipsum}
\usepackage[T1]{fontenc}
\usepackage{icomma}
\usepackage{siunitx}
\usepackage{ragged2e}
\usepackage{comment}

\usepackage{multicol}

\usepackage{esint}

\usepackage{caption}
\usepackage{float}
\usepackage{titling}
\usepackage{blindtext}
\usepackage[square,sort,comma,numbers]{natbib}
\usepackage[colorinlistoftodos]{todonotes}
\definecolor{darkgreen}{rgb}{0.0, 0.4, 0.0}

\begin{comment}
\usepackage{draftwatermark}
\SetWatermarkLightness{0.9}
\DraftwatermarkOptions{color={[rgb]{1,0,0}}}
\SetWatermarkText{REEEEEEE}
\SetWatermarkScale{3}
\end{comment}


%\usepackage[bottom]{footmisc} 
\renewcommand{\footnoterule}{\vfill\kern -3pt \hrule width 0.4\columnwidth \kern 2.6pt}
% either of these moves footnotes to the bottom of the page.

\usepackage{bussproofs}

\usepackage{epigraph}

% \epigraphsize{\small}% Default
\setlength\epigraphwidth{0.75\textwidth}
\setlength\epigraphrule{0pt}
\renewcommand{\beforeepigraphskip}{-0.5\baselineskip}
\renewcommand{\afterepigraphskip}{2\baselineskip}

%\usepackage{etoolbox}
%\makeatletter
%\patchcmd{\epigraph}{\@epitext{#1}}{\itshape\@epitext{#1}}{}{}
%\makeatother

% formats all epigraphs to be italics

%\usepackage{tocbibind}


\setcounter{MaxMatrixCols}{20}

\usepackage{mathdots}

\usepackage{enumitem}


\newcommand{\catname}[1]{{\normalfont\textbf{#1}}}
\newcommand{\Set}{\catname{Set}}
\newcommand{\FinSet}{\catname{FinSet}}
\newcommand{\Rel}{\catname{Rel}}
\newcommand{\Grp}{\catname{Grp}}
\newcommand{\Met}{\catname{Met}}
\newcommand{\Top}{\catname{Top}}
\newcommand{\hTop}{\catname{hTop}}
\newcommand{\Cat}{\catname{Cat}}
\newcommand{\CAT}{\catname{CAT}}
\newcommand{\Mon}{\catname{Mon}}
\newcommand{\Ring}{\catname{Ring}}
\newcommand{\Vect}{\catname{Vect}}
\newcommand{\FinVect}{\catname{FinVect}}
\newcommand{\Mat}{\catname{Mat}}
\newcommand{\Mea}{\catname{Mea}}
\newcommand{\Man}{\catname{Man}}
\newcommand{\Topoi}{\catname{Topoi}}
\newcommand{\Sh}{\catname{Sh}}

\newcommand{\ob}{\operatorname{ob}}
\newcommand{\op}{^{\operatorname{op}}}
\newcommand{\ab}{^{\operatorname{ab}}}
\newcommand{\id}{\operatorname{id}}
\newcommand{\Hom}{\operatorname{Hom}}
\newcommand{\dom}{\operatorname{dom}}
\newcommand{\cod}{\operatorname{cod}}
\newcommand{\ev}{\operatorname{ev}}
\newcommand{\Cone}{\operatorname{Cone}}
\newcommand{\Sub}{\operatorname{Sub}}
\newcommand{\Free}{\operatorname{Free}}


\newcommand{\quotient}[2]{\raisebox{0.1em}{$#1$}\big/\raisebox{-0.1em}{$#2$}}

\DeclareMathOperator*{\colim}{colim}

\DeclareFontFamily{U}{min}{}
\DeclareFontShape{U}{min}{m}{n}{<-> dmjhira}{}



\DeclareMathAlphabet{\mathbbn}{U}{BOONDOX-ds}{m}{n}

\usepackage{adjustbox}
\tikzcdset{scale cd/.style={every label/.append style={scale=#1},
    cells={nodes={scale=#1}}}}



\newlength{\currentparskip}



\graphicspath{ {figures/} }

\addto\captionsenglish{%
    \renewcommand{\listfigurename}{List of Figures}
    \renewcommand{\listtablename}{List of Tables}
}

\usepackage[scr]{rsfso} % this is for Laplace transforms



\usepackage{colortbl}




\usepackage{calrsfs}
\DeclareMathAlphabet{\pazocal}{OMS}{zplm}{m}{n}



\newcommand\scalemath[2]{\scalebox{#1}{\mbox{\ensuremath{\displaystyle #2}}}}







\newcommand\mathitem[1][]{\item[#1]\leavevmode\vspace*{-\dimexpr\baselineskip+\abovedisplayskip\relax}}
% for inserting equations that don't have leading text in itemize or enumerate environments - removes the leading whitespace.


\pagestyle{fancy}
%\renewcommand{\chaptername}{}
%\renewcommand{\chaptermark}[1]{\markboth{#1}{}}
\renewcommand{\sectionmark}[1]{\markright{#1}{}}
%\renewcommand{\subsectionmark}[1]{\markright{#1}{}}


\fancypagestyle{toc}{%
    \fancyhf{}
    \setlength\headheight{15pt}
    \fancyhead[L]{Front Matter}
    \fancyhead[R]{Table of Contents}
    \fancyfoot[C]{Turan Densities | \thepage}
}
\fancypagestyle{lot}{%
    \fancyhf{}
    \setlength\headheight{15pt}
    \fancyhead[L]{Front Matter}
    \fancyhead[R]{List of Tables}
    \fancyfoot[C]{Turan Densities | \thepage}
}
\fancypagestyle{lof}{%
    \fancyhf{}
    \setlength\headheight{15pt}
    \fancyhead[L]{Front Matter}
    \fancyhead[R]{List of Figures}
    \fancyfoot[C]{Turan Densities | \thepage}
}
\fancypagestyle{plain}{%
    \fancyhf{}%
    \renewcommand{\headrulewidth}{0pt}% 
}
\fancypagestyle{preface}{%
    \fancyhf{}%
    \setlength\headheight{15pt}
    \fancyhead[L]{Front Matter}
    \fancyhead[R]{Preface}
    \fancyfoot[C]{Turan Densities | \thepage}
}
\fancypagestyle{main}{%
    \fancyhf{}%
    \setlength\headheight{15pt}
    \fancyhead[L]{REE} %>>>
    \fancyhead[R]{\rightmark}
    \fancyfoot[C]{Turan Densities | \thepage} %>>>
}





\newcommand{\amogus}[1]{\scalebox{#1}{
    \begin{tikzpicture}
        \filldraw[red]
            {[rounded corners=15pt]
            (0,5)
         -- (-1,5)
         -- (-1,1.5)}
         -- (0,1.5)
         -- cycle;
        \begin{scope}
            \clip {[rounded corners=15pt]
                (0,5)
             -- (-1,5)
             -- (-1,1.5)}
             -- (0,1.5)
             -- cycle;
            \filldraw[draw=none, fill=purple, opacity=0.95]
                {[rounded corners=5pt]
                (-1,4.1)
             -- (-0.5,4)}
             -- (0,4.1)
             -- (0,1.5)
             -- (-1,1.5)
             -- cycle;
        \end{scope}
        \draw[line width=4mm]
            {[rounded corners=15pt]
            (0,5)
         -- (-1,5)
         -- (-1,1.5)}
         -- (0,1.5);
         
        \filldraw[red]
            {[rounded corners=10pt]
            (0,3)
         -- (0,0)
         -- (1.6,0)
         -- (1.6,1.5)} 
         -- (3,1.5)
            {[rounded corners=10pt]
         -- (2.5,1.5)
         -- (2.5,0.25)
         -- (4,0.25)}
            {[rounded corners=50pt]
         -- (4,7) 
         -- (0,7)}
         -- cycle;
        \begin{scope}
            \clip
                {[rounded corners=10pt]
                (0,3)
             -- (0,0)
             -- (1.6,0)
             -- (1.6,1.5)} 
             -- (3,1.5)
                {[rounded corners=10pt]
             -- (2.5,1.5)
             -- (2.5,0.25)
             -- (4,0.25)}
                {[rounded corners=50pt]
             -- (4,7) 
             -- (0,7)}
             -- cycle;
            \filldraw[draw=none, fill=purple, opacity=0.95]
                (0,9)
             -- (1,7)
                {[rounded corners=40pt]
             -- (0.4,6)
             -- (0.8,2)
             -- (4,2)}
             -- (4,4.5)
             -- (4,0)
             -- (0,0)
             -- cycle;
        \end{scope}
        \draw[line width=4mm]
            {[rounded corners=15pt]
            (0,3)
         -- (0.15,0)
         -- (1.7,0)} 
            {[rounded corners=5pt]
         -- (1.75,1.5)
         -- (3,1.5)}
         -- (2.35,1.5)
            {[rounded corners=15pt]
         -- (2.4,0.25)
         -- (3.9,0.25)}
            {[rounded corners=15pt]
         -- (4,1.4)
         -- (4.2,2.5)}
            {[rounded corners=45pt]
         -- (3.9,6.8) 
         -- (0.15,7)
         -- cycle};
        \begin{scope}[on background layer]
            \filldraw[draw=none, fill=black, opacity=0.65]
                (1.9,0.1) ellipse (2.9 and 0.7);
        \end{scope}
        \draw[rounded corners=25pt, fill=cyan!40!white, draw=black]
            (4.5,6) rectangle (1.5,4);
        \begin{scope}
            \clip (4.5,6) [rounded corners=25pt] rectangle (1.5,4);
            \filldraw[draw=none, fill=blue!50!black, opacity=0.5]
                {[rounded corners=20pt]
                (2.1,6.2)
             -- (2.1,4.7)
             -- (4.6,4.7)}
             -- (4.6,5.2)
             -- (4.6,4)
             -- (1.5,4)
             -- (1.5,6)
            {[rounded corners=20pt]
             -- (2.5,6)
             -- cycle};
            \draw[very thick, rounded corners=7.5pt, fill=white, draw=none]
                (4,5.5) rectangle (2.8,5);
        \end{scope}
        \draw[rounded corners=25pt, line width=4mm]
            (4.5,5.9) rectangle (1.5,4);
    \end{tikzpicture}
    }
}


\fancypagestyle{sussy}{%
    \fancyhf{}%
    \setlength\headheight{15pt}
    \fancyhead[L]{\amogus{0.05}} %>>>
    \fancyhead[R]{\rightmark}
    \fancyfoot[C]{Turan Densities | \thepage} %>>>
}
